\section{Limitations}
\label{sec:limitations}

\emph{One budget.} Every conclusion concerns eight backward-looking price features and one gradient-boosting
configuration. The ablations show that the conclusions are stable to the decision threshold and to a much smaller
model, but richer information---fundamentals, order flow, cross-sectional or alternative data---could change them.
MBCA is the protocol for testing that fairly, not a claim that the space of learned components is exhausted.

\emph{Daily closes only.} Signals, stops, and fills use daily closes, so intraday stop hits are not modelled. The
close-fill convention is conservative relative to stop-level fills, but overnight gaps and liquidity at the open are
not simulated.

\emph{Data.} Study~2 relies on one primary vendor, cross-checked against FRED and MarketWatch. Continuous futures are
unadjusted front-month rolls; single stocks are current large capitalizations and the crypto series are surviving
coins. These survivorship-prone markets are analysed separately and do not drive the conclusions, and one planned
symbol was unavailable.

\emph{Portfolio construction.} Equal weighting across instruments with daily mark-to-market accounting is transparent
but simple. Volatility targeting or risk parity would change the absolute numbers, though not the paired
variant-versus-control design.

\emph{Statistics.} The block bootstrap assumes stationarity within each window, and FDR control is applied within
pre-defined families rather than across the whole study. Most market-level effects are fragile to dropping a single
instrument; only the significant ones are shown to be robust to it.

\emph{Scope.} The evidence covers five rule families on daily data. MBCA itself is domain agnostic, but we have not
applied it outside trading.
